\section{Introduction}
A time-reversal-invariant conductor can carry a transverse current at second order in an electric field. In the anomalous-velocity contribution, the electric field probes Berry curvature through the first-order nonequilibrium occupation. A constant relaxation time reduces this occupation to velocity times a lifetime, leading to the Berry-curvature-dipole expression~\cite{SodemannFu2015}. The geometric weight and the kinetic distribution then appear almost separable. Microscopic multivalley collisions remove that simplification: carriers can change valley as well as direction, and distinct distortions relax at different rates.

This distinction matters when reducing a multivalley description. Suppose an occupied valley has zero Berry curvature everywhere on its conducting contour. It contributes no direct weight to the anomalous-velocity Hall functional. It nevertheless carries longitudinal current, accepts carriers from other valleys and returns a driven distribution to them. Its absence from the Hall readout does not imply its absence from the kinetic equation. We call such a sector \emph{Hall silent}; the term refers to a specified response channel, not to zero conductivity or to an inert band.

Scattering-induced nonlinear Hall redistribution is already established. Lihm and Park identified valley-polarizing relaxation modes and developed scattering-channel diagnostics, including an exact response difference after channel subtraction~\cite{LihmPark2024}. Multichannel kinetic theories also show why the anomalous-velocity contribution alone need not exhaust the nonlinear Hall response~\cite{NandySodemann2019,Du2019}. We address the narrower question of when an entire conducting, Hall-silent sector may be omitted at a stated tolerance. This differs from asking whether a single relaxation time describes all valleys or whether a particular scattering channel enhances a response.

Our central result is that a zero-curvature conducting valley can remain kinetically indispensable. Safe omission depends on the requested observable, the active and passive relaxation structure, and the accepted error. Within an explicit four-valley model, a contour-weighted projection separates scalar rescaling from Hall-sensitive redistribution. A fixed physical hierarchy and its invariant-subspace residual explain why either contribution can dominate: the relevant competition is between the change in a reference-current amplitude and the excitation of an orthogonal valley distortion with a different Hall overlap and relaxation stiffness. Opposing contributions can also make safe omission cancellation-dependent under a fixed-budget test. We connect that statement to 1\%, 5\% and 10\% decisions, rather than inferring practical significance from a small cancellation factor alone.

The contribution is the use of sector elimination and observable-error analysis in a multivalley omission problem, together with a physically interpretable hierarchy identifying which transport degrees of freedom control the error. The mathematical tools are standard. We use Schur elimination, closely related to Kron reduction of network operators~\cite{DorflerBullo2013}, and an observable-error identity of the adjoint type~\cite{PierceGiles2004}. Their application exposes physical omission mechanisms; no mathematical priority claim is made. A small passive mass supplies a further control: it converts a pre-existing passive valley displacement into direct Hall readout, even without intersector coupling. Finally, ordinary conductivity makes omission of the same sector a different question for an open-circuit signal.

All results concern the anomalous-velocity nonlinear Hall contribution of the specified homogeneous semiclassical model. They do not give a quantitative material prediction, a universal omission threshold, or the complete nonlinear Hall conductivity. The Supplement contains the detailed numerical validation, disorder controls, frequency analysis and sufficient error bounds; the main text develops the physical explanation.
